\section{Data and Methods}
\subsection{Data and benchmark}
The shared snapshot contains 2,185 usable Copenhagen delivery days from
2019--2024. Inputs include one- and seven-day price lags for DK1, DK2,
DE-LU, NL, and SE4 \cite{energinet,energycharts}, calendar features, and
archived JMA forecasts of temperature, wind, and cloud cover
\cite{openmeteo}. The intended issue time is 10:00 Copenhagen time on the
preceding day. Historical release times and revisions are partly audited.
Daily training curves interpolate the missing spring hour and average the
repeated autumn hour. Scoring matches UTC timestamps, retaining the 23
or 25 actual hours on transition days.

The calendar and price-lag features represent daily and weekly
seasonality and link DK1 to neighbouring markets. For weather, the
Open-Meteo \texttt{previous\_day2} archive provides forecasts for a
fixed lead. Using a forecast archive avoids substituting realized
weather for forecasts, while the archive lead leaves the availability
of particular historical versions at issue time uncertain.

We refit each model monthly using all available earlier days. We report
MAE and RMSE in EUR/MWh and count predictions within \(\pm5\%\) of the
actual price among hours with \(|\mathrm{actual}|\geq20\) EUR/MWh.
Ridge uses standardized inputs and fixed \(\alpha=100\).

\subsection{Forecast models}
Lasso selects \(\alpha=0.3\) by RMSE over 2023 rolling validation months,
with training-window scaling and a 10,000-iteration limit. The reported
Random Forest uses 100 trees and a minimum leaf size of five. The
2022 MAE search selected two, which we report as a sensitivity check.
Histogram Gradient Boosting fits one estimator per clock hour
\cite{friedman2001,pedregosa2011}. The 60-tree starter
(learning rate 0.08, 15 leaves, \(L_2=1\)) led the three tested variants
on 2023 MAE and hit count. The MLP predicts all 24 hours jointly
\cite{lago2021,goodfellow2016}. Rolling 2023 validation selected a
40-unit tanh layer with \(L_2\) penalty \(\alpha=10\). Its input and
target scalers are fitted within each training window.

The model families span different ways to use the same information.
Ridge shrinks coefficients when lagged prices and weather inputs are
correlated, while Lasso can set coefficients to zero. The tree models
can represent nonlinear interactions, and the forest and MLP predict
the 24-hour curve jointly. We disabled automatic early stopping for
the MLP to avoid a random validation split, and compared layer width,
penalty, and iteration budget in chronological folds.

\subsection{K-means regimes}
Each day is described by the previous DK1 price mean, standard deviation,
and evening-minus-morning spread, plus the mean forecast temperature,
wind speed, and cloud cover. Scaling balances units in Euclidean distance.
We fit the scaler and K-means on 1,454 pre-2023 days and freeze them for
2023--2024 \cite{clustering}. Across \(k=2,\ldots,8\), we considered the
Davies--Bouldin index, smallest cluster size, and interpretability.
We kept \(k=4\), whose smallest training cluster was 12.9\% of days,
compared with 6.2\% for \(k=5\).

\subsection{Apriori and rule transfer}
An hourly basket contains season, six-hour time band, and market regime.
We use Apriori \cite{apriori} to mine antecedents supported by at least
5\% of 2023 hours and matching at least 20 large Ridge errors.
A large error exceeds the 2023 absolute-error 90th percentile.
Confidence is the large-error rate among matching hours. Lift divides
confidence by the annual large-error rate. We retain the highest-lift
calendar-only and cluster-containing antecedents, breaking ties by
support and item order. We apply both antecedents unchanged in 2024,
using each model's fixed 2023 threshold. We also check a common Ridge
threshold for all models.

For the cluster rule, we compare cluster 0 with other regimes within
12:00--17:59. As an out-of-course method, a moving-block bootstrap
resamples contiguous seven-day blocks of 2024 delivery days
(5,000 draws). Keeping all hours in each day together preserves
short-range dependence. The percentile interval describes uncertainty
conditional on the selected rule and threshold.
The resampling computes the rate difference for the two afternoon
groups in each draw, using the same sampled days for both. This keeps
dependence among hours on the same day and among neighbouring days
within a block, which would be lost with independent hourly draws.
